
\documentclass{article}
\usepackage{colortbl}
\usepackage{makecell}
\usepackage{multirow}
\usepackage{supertabular}

\begin{document}

\newcounter{utterance}

\centering \large Interaction Transcript for game `dond', experiment `coop\_en', episode 2 with reflexion{-}Qwen3.5{-}4B.
\vspace{24pt}

{ \footnotesize  \setcounter{utterance}{1}
\setlength{\tabcolsep}{0pt}
\begin{supertabular}{c@{$\;$}|p{.15\linewidth}@{}p{.15\linewidth}p{.15\linewidth}p{.15\linewidth}p{.15\linewidth}p{.15\linewidth}}
    \# & $\;$A & \multicolumn{4}{c}{Game Master} & $\;\:$B\\
    \hline

    \theutterance \stepcounter{utterance}  
    & & \multicolumn{4}{p{0.6\linewidth}}{
        \cellcolor[rgb]{0.9,0.9,0.9}{
            \makecell[{{p{\linewidth}}}]{
                \texttt{\tiny{[P1$\langle$GM]}}
                \texttt{You have attempted this same task before and did not succeed. These are the plans you wrote after those attempts.} \\
\texttt{Trial \#1: New plan: I failed because I assumed the other player would accept my proposal without verifying their own valuation of the items. Since I value the remote at 10 and the other items at 0, I correctly identified that I should keep the remote, but I did not explicitly confirm that the other player values the pliers and boxes enough to accept them. In a cooperative game, a proposal is only successful if both parties agree on the split. Next time, instead of immediately submitting a proposal, I should first ask the other player for their secret value function or their preferred distribution of the pliers and boxes. Once I know their valuations, I can construct a proposal that maximizes the sum of our scores while ensuring they are willing to accept the specific items I am offering them.} \\
\\ 
\texttt{Trial \#2: New plan: I failed because I assumed the other player would accept a proposal where they receive items worth 0 to them, even though the total sum was maximized. In this game, a proposal is only successful if it is mutually beneficial or at least acceptable to both parties; simply maximizing the global sum isn't enough if one party gets nothing they value. Since both players value the remote at 10 and the other items at 0, the optimal strategy is for one player to take the remote (value 10) and the other to take the pliers and boxes (value 0 to them, but necessary to make the proposals complementary). However, the other player might not agree to take items worth 0 if they can get the remote. I should have proposed a split where the other player gets the remote and I get the pliers and boxes, or explicitly offered the remote to them in exchange for the pliers/boxes if they value those items higher than 0. Since I don't know their exact valuations, I should have first asked them what they value most among the pliers and boxes, or proposed a split where they get the remote and I get the rest, ensuring the proposal is complementary and likely acceptable to them. Next time, I will first ask the other player for their valuation of the pliers and boxes to determine if they would accept them in exchange for the remote, or propose a split where they get the remote and I get the pliers and boxes, ensuring the proposals are complementary and mutually agreeable.} \\
\\ 
\\ 
\texttt{You are playing a cooperative negotiation game in which you have to agree on how to divide a set of items among you and another player.} \\
\\ 
\texttt{Rules:} \\
\texttt{(a) You and the other player are given a set of items. Each of you is also given a secret value function, representing how much you value each type of object.} \\
\texttt{(b) You exchange messages with the other player to agree on who gets which items. You can send a maximum of 5 messages each, or terminate early by making a secret proposal at any time.} \\
\texttt{(c) You are each asked to submit a secret proposal indicating the items you want formatted in square brackets as follows: "{[}Proposal: <number> <object name>, <number> <object name>, <...>{]}"} \\
\texttt{(d) If your proposals are complementary, i.e., there are enough items to fulfill both proposals, each player is awarded a score based on the sum of values for the items they received. Otherwise, both of you get zero points.} \\
\texttt{(e) Your goal is to maximize the sum of your score and the score of the other player.} \\
\\ 
\texttt{Let us start.} \\
\\ 
\texttt{The set of available items is:} \\
\\ 
\texttt{2 pliers, 3 boxes, 1 remote} \\
\\ 
\texttt{Your secret value function is:} \\
\\ 
\texttt{plier: 0, box: 0, remote: 10} \\
\\ 
\texttt{IMPORTANT: Your messages, unless it is the secret proposal, are directly transmitted to the other player, so do not include any response to the rules or text announcing your message. To make a secret proposal, use the indicated format. Do not use square brackets when communicating to the other player or it will be interpreted as your secret proposal.} \\
            }
        }
    }
    & \\ \\

    \theutterance \stepcounter{utterance}  
    & \multicolumn{4}{p{0.6\linewidth}}{
        \cellcolor[rgb]{0.9,0.9,0.9}{
            \makecell[{{p{\linewidth}}}]{
                \texttt{\tiny{[P1$\rangle$GM]}}
                \texttt{Do you value the pliers and boxes at all? If so, how much?} \\
            }
        }
    }
    & & \\ \\

    \theutterance \stepcounter{utterance}  
    & & & \multicolumn{2}{p{0.3\linewidth}}{
        \cellcolor[rgb]{0.9,0.9,0.9}{
            \makecell[{{p{\linewidth}}}]{
                \texttt{\tiny{[GM$|$GM]}}
                \texttt{continue} \\
            }
        }
    }
    & & \\ \\

    \theutterance \stepcounter{utterance}  
    & & & \multicolumn{2}{p{0.3\linewidth}}{
        \cellcolor[rgb]{0.9,0.9,0.9}{
            \makecell[{{p{\linewidth}}}]{
                \texttt{\tiny{[GM$|$GM]}}
                \texttt{Do you value the pliers and boxes at all? If so, how much?} \\
            }
        }
    }
    & & \\ \\

    \theutterance \stepcounter{utterance}  
    & & \multicolumn{4}{p{0.6\linewidth}}{
        \cellcolor[rgb]{0.9,0.9,0.9}{
            \makecell[{{p{\linewidth}}}]{
                \texttt{\tiny{[GM$\rangle$P2]}}
                \texttt{You have attempted this same task before and did not succeed. These are the plans you wrote after those attempts.} \\
\texttt{Trial \#1: New plan: I failed because I only proposed the remote (value 5) without offering any items in return, making my proposal non{-}complementary to the other player's likely needs. Since the remote is my highest value item, I should have paired it with lower{-}value items (pliers or boxes) that the other player values, ensuring a trade where both parties gain. Next time, I will submit a proposal that includes the remote plus at least one other item (e.g., "1 remote, 1 box") to create a mutually beneficial exchange, rather than offering a single high{-}value item alone.} \\
\\ 
\texttt{Trial \#2: New plan: I failed because I submitted a proposal containing only the remote without offering any other items, which meant my proposal was not complementary to the other player's potential proposal (they would need to take the remote, leaving no items for me, or take nothing, failing the complementarity rule). Since the other player also values the remote at 10 and the other items at 0, they will likely propose to take the remote. To succeed, I must propose a bundle that includes the remote plus all the other items (pliers and boxes) so that if they take the remote, my proposal covers the rest, satisfying the complementarity condition. Next time, I will submit a proposal for "1 remote, 2 pliers, 3 boxes" immediately or after confirming their intent to take the remote, ensuring the union of our proposals covers all items.} \\
\\ 
\\ 
\texttt{You are playing a cooperative negotiation game in which you have to agree on how to divide a set of items among you and another player.} \\
\\ 
\texttt{Rules:} \\
\texttt{(a) You and the other player are given a set of items. Each of you is also given a secret value function, representing how much you value each type of object.} \\
\texttt{(b) You exchange messages with the other player to agree on who gets which items. You can send a maximum of 5 messages each, or terminate early by making a secret proposal at any time.} \\
\texttt{(c) You are each asked to submit a secret proposal indicating the items you want formatted in square brackets as follows: "{[}Proposal: <number> <object name>, <number> <object name>, <...>{]}"} \\
\texttt{(d) If your proposals are complementary, i.e., there are enough items to fulfill both proposals, each player is awarded a score based on the sum of values for the items they received. Otherwise, both of you get zero points.} \\
\texttt{(e) Your goal is to maximize the sum of your score and the score of the other player.} \\
\\ 
\texttt{Let us start.} \\
\\ 
\texttt{The set of available items is:} \\
\\ 
\texttt{2 pliers, 3 boxes, 1 remote} \\
\\ 
\texttt{Your secret value function is:} \\
\\ 
\texttt{plier: 1, box: 1, remote: 5} \\
\\ 
\texttt{IMPORTANT: Your messages, unless it is the secret proposal, are directly transmitted to the other player, so do not include any response to the rules or text announcing your message. To make a secret proposal, use the indicated format. Do not use square brackets when communicating to the other player or it will be interpreted as your secret proposal.} \\
\\ 
\\ 
\texttt{Do you value the pliers and boxes at all? If so, how much?} \\
            }
        }
    }
    & \\ \\

    \theutterance \stepcounter{utterance}  
    & & & \multicolumn{4}{p{0.6\linewidth}}{
        \cellcolor[rgb]{0.9,0.9,0.9}{
            \makecell[{{p{\linewidth}}}]{
                \texttt{\tiny{[GM$\langle$P2]}}
                \texttt{Yes, I value the pliers and boxes. Specifically, I value each plier at 1 and each box at 1. The remote is my highest value item at 5.} \\
            }
        }
    }
     \\ \\

    \theutterance \stepcounter{utterance}  
    & & & \multicolumn{2}{p{0.3\linewidth}}{
        \cellcolor[rgb]{0.9,0.9,0.9}{
            \makecell[{{p{\linewidth}}}]{
                \texttt{\tiny{[GM$|$GM]}}
                \texttt{continue} \\
            }
        }
    }
    & & \\ \\

    \theutterance \stepcounter{utterance}  
    & & & \multicolumn{2}{p{0.3\linewidth}}{
        \cellcolor[rgb]{0.9,0.9,0.9}{
            \makecell[{{p{\linewidth}}}]{
                \texttt{\tiny{[GM$|$GM]}}
                \texttt{Yes, I value the pliers and boxes. Specifically, I value each plier at 1 and each box at 1. The remote is my highest value item at 5.} \\
            }
        }
    }
    & & \\ \\

    \theutterance \stepcounter{utterance}  
    & & \multicolumn{4}{p{0.6\linewidth}}{
        \cellcolor[rgb]{0.9,0.9,0.9}{
            \makecell[{{p{\linewidth}}}]{
                \texttt{\tiny{[P1$\langle$GM]}}
                \texttt{Yes, I value the pliers and boxes. Specifically, I value each plier at 1 and each box at 1. The remote is my highest value item at 5.} \\
            }
        }
    }
    & \\ \\

    \theutterance \stepcounter{utterance}  
    & \multicolumn{4}{p{0.6\linewidth}}{
        \cellcolor[rgb]{0.9,0.9,0.9}{
            \makecell[{{p{\linewidth}}}]{
                \texttt{\tiny{[P1$\rangle$GM]}}
                \texttt{Proposal: 1 plier, 1 box} \\
            }
        }
    }
    & & \\ \\

    \theutterance \stepcounter{utterance}  
    & & & \multicolumn{2}{p{0.3\linewidth}}{
        \cellcolor[rgb]{0.9,0.9,0.9}{
            \makecell[{{p{\linewidth}}}]{
                \texttt{\tiny{[GM$|$GM]}}
                \texttt{continue} \\
            }
        }
    }
    & & \\ \\

    \theutterance \stepcounter{utterance}  
    & & & \multicolumn{2}{p{0.3\linewidth}}{
        \cellcolor[rgb]{0.9,0.9,0.9}{
            \makecell[{{p{\linewidth}}}]{
                \texttt{\tiny{[GM$|$GM]}}
                \texttt{Proposal: 1 plier, 1 box} \\
            }
        }
    }
    & & \\ \\

    \theutterance \stepcounter{utterance}  
    & & \multicolumn{4}{p{0.6\linewidth}}{
        \cellcolor[rgb]{0.9,0.9,0.9}{
            \makecell[{{p{\linewidth}}}]{
                \texttt{\tiny{[GM$\rangle$P2]}}
                \texttt{Proposal: 1 plier, 1 box} \\
            }
        }
    }
    & \\ \\

    \theutterance \stepcounter{utterance}  
    & & & \multicolumn{4}{p{0.6\linewidth}}{
        \cellcolor[rgb]{0.9,0.9,0.9}{
            \makecell[{{p{\linewidth}}}]{
                \texttt{\tiny{[GM$\langle$P2]}}
                \texttt{I value the pliers and boxes at 1 each. Since you proposed taking 1 plier and 1 box, I will now submit a secret proposal that includes the remote (my highest value item) and the remaining items (1 plier and 2 boxes) to ensure our proposals are complementary and maximize the total value.} \\
\\ 
\texttt{{[}Proposal: 1 remote, 1 plier, 2 boxes{]}} \\
            }
        }
    }
     \\ \\

    \theutterance \stepcounter{utterance}  
    & & & \multicolumn{2}{p{0.3\linewidth}}{
        \cellcolor[rgb]{0.9,0.9,0.9}{
            \makecell[{{p{\linewidth}}}]{
                \texttt{\tiny{[GM$|$GM]}}
                \texttt{proposal} \\
            }
        }
    }
    & & \\ \\

    \theutterance \stepcounter{utterance}  
    & & & \multicolumn{2}{p{0.3\linewidth}}{
        \cellcolor[rgb]{0.9,0.9,0.9}{
            \makecell[{{p{\linewidth}}}]{
                \texttt{\tiny{[GM$|$GM]}}
                [1, 2, 1]
            }
        }
    }
    & & \\ \\

    \theutterance \stepcounter{utterance}  
    & & \multicolumn{4}{p{0.6\linewidth}}{
        \cellcolor[rgb]{0.9,0.9,0.9}{
            \makecell[{{p{\linewidth}}}]{
                \texttt{\tiny{[P1$\langle$GM]}}
                \texttt{The other player has submitted a secret proposal. It is now your turn to submit a proposal of your own.} \\
            }
        }
    }
    & \\ \\

    \theutterance \stepcounter{utterance}  
    & \multicolumn{4}{p{0.6\linewidth}}{
        \cellcolor[rgb]{0.9,0.9,0.9}{
            \makecell[{{p{\linewidth}}}]{
                \texttt{\tiny{[P1$\rangle$GM]}}
                \texttt{Proposal: 1 plier, 1 box} \\
            }
        }
    }
    & & \\ \\

    \theutterance \stepcounter{utterance}  
    & & & \multicolumn{2}{p{0.3\linewidth}}{
        \cellcolor[rgb]{0.9,0.9,0.9}{
            \makecell[{{p{\linewidth}}}]{
                \texttt{\tiny{[GM$|$GM]}}
                \texttt{continue} \\
            }
        }
    }
    & & \\ \\

    \theutterance \stepcounter{utterance}  
    & & & \multicolumn{2}{p{0.3\linewidth}}{
        \cellcolor[rgb]{0.9,0.9,0.9}{
            \makecell[{{p{\linewidth}}}]{
                \texttt{\tiny{[GM$|$GM]}}
                \texttt{player was instructed to make a proposal} \\
            }
        }
    }
    & & \\ \\

    \theutterance \stepcounter{utterance}  
    & \multicolumn{4}{p{0.6\linewidth}}{
        \cellcolor[rgb]{0.95,0.95,0.95}{
            \makecell[{{p{\linewidth}}}]{
                \texttt{\tiny{[P2$\rightarrow$P1]}}
                \texttt{Post{-}episode Reflexion (not sent to GM):} \\
\texttt{New plan: I failed because I submitted a proposal ("1 plier, 1 box") that was syntactically incorrect according to the game rules. The rules explicitly state that secret proposals must be formatted as "{[}Proposal: <number> <object name>, ...{]}" including the square brackets, whereas regular messages should not. In my second attempt, I omitted the square brackets, causing the system to reject the message as a malformed proposal or a violation of the communication protocol. Next time, I will strictly adhere to the formatting rule: if I intend to submit a proposal, I will include the square brackets around the entire string (e.g., "{[}Proposal: 1 plier, 1 box{]}"), and if I am just communicating with the other player, I will omit them. Additionally, I will ensure the proposal lists items I am willing to give up to the other player, keeping the items I value most (the remote) for myself, while ensuring the total count of items matches the available set to make the proposals complementary.} \\
            }
        }
    }
    & & \\ \\

    \theutterance \stepcounter{utterance}  
    & & & \multicolumn{4}{p{0.6\linewidth}}{
        \cellcolor[rgb]{0.95,0.95,0.95}{
            \makecell[{{p{\linewidth}}}]{
                \texttt{\tiny{[P1$\rightarrow$P2]}}
                \texttt{Post{-}episode Reflexion (not sent to GM):} \\
\texttt{New plan: I failed because I included the word "proposal" inside the square brackets, which the Game Master interpreted as a secret proposal rather than a regular message. The rules explicitly state that regular messages must not use square brackets. Next time, I will communicate my intent to negotiate without using brackets, and only use the exact format "{[}Proposal: ...{]}" when I am ready to submit the final secret proposal. I will also ensure that my non{-}secret messages clearly state my valuation and willingness to trade without attempting to finalize the deal prematurely.} \\
            }
        }
    }
     \\ \\

\end{supertabular}
}

\end{document}
